% --- Archon physics formalization source begin ---
% archon:physics
% archon:covers PhyXMiniProblems/problem_phyx_mini_0247.lean
% archon:source-report reports/phyx_mini/problem_phyx_mini_0247.source.json

\chapter{Physics problem phyx\_mini\_0247}
\label{ch:PhyXMiniProblems_problem_phyx_mini_0247}

\paragraph{Problem source.}
\#\# Physical scenario

Your firm has been hired to design a system that allows airplane pilots to make instrument landings in rain or fog. You've decided to place two radio transmitters \$50 \textbackslash{}, \textbackslash{}mathrm\{m\}\$ apart on either side of the runway. These two transmitters will broadcast the same frequency, but out of phase with each other. This will cause a nodal line to extend straight off the end of the runway. As long as the airplane's receiver is silent, the pilot knows she's directly in line with the runway. If she drifts to one side or the other, the radio will pick up a signal and sound a warning beep. To have sufficient accuracy, the first intensity maxima need to be \$60 \textbackslash{}, \textbackslash{}mathrm\{m\}\$ on either side of the nodal line at a distance of \$3.0 \textbackslash{}, \textbackslash{}mathrm\{km\}\$.

\#\# Question

What frequency should you specify for the transmitters?

\#\# Answer choices

- A: A: \$150\textbackslash{};MHz\$

- B: B: \$100\textbackslash{};MHz\$

- C: C: \$200\textbackslash{};MHz\$

- D: D: \$120\textbackslash{};MHz\$

\#\# Figure caption (auxiliary; use the image as primary evidence)

The illustration is a schematic representation of wave interference patterns with the following elements:

1. **Wave Sources:** Two identical sources, represented by stars within circles, are shown at the left. Circular wave fronts emanate from these sources.

2. **Runway:** A rectangular shaded area labeled as "Runway" measures 50 meters wide.

3. **Interference Patterns:** 
   - **Antinodal lines** (labeled) extend from the sources, indicating regions of constructive interference.
   - **Nodal line** (labeled) is also shown, representing a region of destructive interference between the wave sources.

4. **Lines and Points:**
   - Lines \textbackslash{}( r\_1 \textbackslash{}) and \textbackslash{}( r\_2 \textbackslash{}) extend from the wave sources to a point labeled \textbackslash{}( P \textbackslash{}) on the upper right, indicating paths of interference.
   - Point \textbackslash{}( P \textbackslash{}) is located 60 meters above the nodal line.
   - A horizontal line extends 3000 meters to the right from the runway area.

5. **Phase Difference:**
   - A note near the bottom left indicates a phase difference: \textbackslash{}(\textbackslash{}Delta \textbackslash{}phi\_0 = \textbackslash{}pi \textbackslash{}, \textbackslash{}text\{rad\}\textbackslash{}).

The diagram illustrates how wave interference occurs with nodal and antinodal lines indicating destructive and constructive interference, respectively.

\#\# Dataset metadata

Wave Properties; Spatial Relation Reasoning, Multi-Formula Reasoning

\paragraph{Recorded answer/context.}
A: A: \$150\textbackslash{};MHz\$

\paragraph{Figure/image path.}
/root/proposal\_for\_physic/science-mango/phyx\_mini\_run/phyx\_data/test\_image/247.png

\paragraph{Formalization target.}
create a compiling Lean file with sorry bodies at `PhyXMiniProblems/problem\_phyx\_mini\_0247.lean`. The Lean declarations must preserve the physical quantities, dimensions or dimensional roles, figure labels, governing-law hypotheses, and final relation expressed by this problem.
Use Mathlib/Physlib names found through LeanExplore where available. If a physics API is missing, introduce faithful local abstractions rather than scalar placeholder aliases.

\begin{theorem}[Physics formalization target]
\label{thm:physics:phyx_mini_0247:target}
\lean{PhyXMiniProblems.ProblemPhyXMini0247.problem_phyx_mini_0247}
\uses{def:physics:phyx-mini-0247:phyxminiproblems-problemphyxmini0247-radiolandinginterferencesetup, def:physics:phyx-mini-0247:phyxminiproblems-problemphyxmini0247-matchescoherentradiosourcedescription, def:physics:phyx-mini-0247:phyxminiproblems-problemphyxmini0247-matchesproblemandfigurereadouts, def:physics:phyx-mini-0247:phyxminiproblems-problemphyxmini0247-matchessuppliedradiolandingfigure, def:physics:phyx-mini-0247:phyxminiproblems-problemphyxmini0247-matcheslandingguidancedescription, def:physics:phyx-mini-0247:phyxminiproblems-problemphyxmini0247-hasphysicalradiointerferenceparameters, def:physics:phyx-mini-0247:phyxminiproblems-problemphyxmini0247-satisfiespathdifferencedefinition, def:physics:phyx-mini-0247:phyxminiproblems-problemphyxmini0247-satisfiescontrolledfarfieldpathdifferencemodel, def:physics:phyx-mini-0247:phyxminiproblems-problemphyxmini0247-satisfiesoutofphaseinterferencecriterion, def:physics:phyx-mini-0247:phyxminiproblems-problemphyxmini0247-satisfiesradiowavespeedlaw, def:physics:phyx-mini-0247:phyxminiproblems-problemphyxmini0247-usesambientairvacuumpropagation, def:physics:phyx-mini-0247:phyxminiproblems-problemphyxmini0247-recordeddatasetanswer, def:physics:phyx-mini-0247:phyxminiproblems-problemphyxmini0247-isuniquematchinganswerchoice}
The assigned autoformalize agent should translate this physics problem into Lean declarations in the covered file, with theorem and lemma proof bodies written as `by sorry`.
\end{theorem}
\begin{proof}
This is an autoformalization task, not a proof task. Produce statements that can later be proved by the physics prover without weakening the physical model.
\end{proof}

% --- Archon declaration topology begin ---
\section{Declaration topology}
\begin{definition}[Length Quantity]
  \label{def:physics:phyx-mini-0247:phyxminiproblems-problemphyxmini0247-lengthquantity}
  \lean{PhyXMiniProblems.ProblemPhyXMini0247.LengthQuantity}
  A nonnegative unit-independent physical length
\end{definition}
\begin{proof}
  This is the declared model definition of Length Quantity.
\end{proof}
\begin{definition}[Frequency Quantity]
  \label{def:physics:phyx-mini-0247:phyxminiproblems-problemphyxmini0247-frequencyquantity}
  \lean{PhyXMiniProblems.ProblemPhyXMini0247.FrequencyQuantity}
  A nonnegative physical frequency, carrying inverse-time dimension
\end{definition}
\begin{proof}
  This is the declared model definition of Frequency Quantity.
\end{proof}
\begin{definition}[Speed Quantity]
  \label{def:physics:phyx-mini-0247:phyxminiproblems-problemphyxmini0247-speedquantity}
  \lean{PhyXMiniProblems.ProblemPhyXMini0247.SpeedQuantity}
  A nonnegative physical propagation speed
\end{definition}
\begin{proof}
  This is the declared model definition of Speed Quantity.
\end{proof}
\begin{definition}[length Readout]
  \label{def:physics:phyx-mini-0247:phyxminiproblems-problemphyxmini0247-lengthreadout}
  \lean{PhyXMiniProblems.ProblemPhyXMini0247.lengthReadout}
  \uses{def:physics:phyx-mini-0247:phyxminiproblems-problemphyxmini0247-lengthquantity}
  Read a physical length in a selected length unit
\end{definition}
\begin{proof}
  This is the declared model definition of length Readout.
\end{proof}
\begin{definition}[frequency Readout]
  \label{def:physics:phyx-mini-0247:phyxminiproblems-problemphyxmini0247-frequencyreadout}
  \lean{PhyXMiniProblems.ProblemPhyXMini0247.frequencyReadout}
  \uses{def:physics:phyx-mini-0247:phyxminiproblems-problemphyxmini0247-frequencyquantity}
  Read a physical frequency in inverse units of a selected time unit
\end{definition}
\begin{proof}
  This is the declared model definition of frequency Readout.
\end{proof}
\begin{definition}[speed Readout]
  \label{def:physics:phyx-mini-0247:phyxminiproblems-problemphyxmini0247-speedreadout}
  \lean{PhyXMiniProblems.ProblemPhyXMini0247.speedReadout}
  \uses{def:physics:phyx-mini-0247:phyxminiproblems-problemphyxmini0247-speedquantity}
  Read a physical speed in selected length and time units
\end{definition}
\begin{proof}
  This is the declared model definition of speed Readout.
\end{proof}
\begin{definition}[length In Meters]
  \label{def:physics:phyx-mini-0247:phyxminiproblems-problemphyxmini0247-lengthinmeters}
  \lean{PhyXMiniProblems.ProblemPhyXMini0247.lengthInMeters}
  \uses{def:physics:phyx-mini-0247:phyxminiproblems-problemphyxmini0247-lengthquantity, def:physics:phyx-mini-0247:phyxminiproblems-problemphyxmini0247-lengthreadout}
  Meter readout used for path and far-field equations
\end{definition}
\begin{proof}
  This is the declared model definition of length In Meters.
\end{proof}
\begin{definition}[length In Kilometers]
  \label{def:physics:phyx-mini-0247:phyxminiproblems-problemphyxmini0247-lengthinkilometers}
  \lean{PhyXMiniProblems.ProblemPhyXMini0247.lengthInKilometers}
  \uses{def:physics:phyx-mini-0247:phyxminiproblems-problemphyxmini0247-lengthquantity, def:physics:phyx-mini-0247:phyxminiproblems-problemphyxmini0247-lengthreadout}
  Kilometer readout used by the stated 3.0 km downrange distance
\end{definition}
\begin{proof}
  This is the declared model definition of length In Kilometers.
\end{proof}
\begin{definition}[frequency In Hertz]
  \label{def:physics:phyx-mini-0247:phyxminiproblems-problemphyxmini0247-frequencyinhertz}
  \lean{PhyXMiniProblems.ProblemPhyXMini0247.frequencyInHertz}
  \uses{def:physics:phyx-mini-0247:phyxminiproblems-problemphyxmini0247-frequencyquantity, def:physics:phyx-mini-0247:phyxminiproblems-problemphyxmini0247-frequencyreadout}
  Hertz readout, i.e. cycles per SI second
\end{definition}
\begin{proof}
  This is the declared model definition of frequency In Hertz.
\end{proof}
\begin{definition}[frequency In Megahertz]
  \label{def:physics:phyx-mini-0247:phyxminiproblems-problemphyxmini0247-frequencyinmegahertz}
  \lean{PhyXMiniProblems.ProblemPhyXMini0247.frequencyInMegahertz}
  \uses{def:physics:phyx-mini-0247:phyxminiproblems-problemphyxmini0247-frequencyquantity, def:physics:phyx-mini-0247:phyxminiproblems-problemphyxmini0247-frequencyinhertz}
  Megahertz readout of a physical frequency
\end{definition}
\begin{proof}
  This is the declared model definition of frequency In Megahertz.
\end{proof}
\begin{definition}[speed In Meters Per Second]
  \label{def:physics:phyx-mini-0247:phyxminiproblems-problemphyxmini0247-speedinmeterspersecond}
  \lean{PhyXMiniProblems.ProblemPhyXMini0247.speedInMetersPerSecond}
  \uses{def:physics:phyx-mini-0247:phyxminiproblems-problemphyxmini0247-speedquantity, def:physics:phyx-mini-0247:phyxminiproblems-problemphyxmini0247-speedreadout}
  Meter-per-second readout of a physical speed
\end{definition}
\begin{proof}
  This is the declared model definition of speed In Meters Per Second.
\end{proof}
\begin{definition}[vacuum Speed Of Light In Meters Per Second]
  \label{def:physics:phyx-mini-0247:phyxminiproblems-problemphyxmini0247-vacuumspeedoflightinmeterspersecond}
  \lean{PhyXMiniProblems.ProblemPhyXMini0247.vacuumSpeedOfLightInMetersPerSecond}
  Physlib's exact vacuum speed-of-light readout, in meters per second
\end{definition}
\begin{proof}
  This is the declared model definition of vacuum Speed Of Light In Meters Per Second.
\end{proof}
\begin{definition}[Transmitter Label]
  \label{def:physics:phyx-mini-0247:phyxminiproblems-problemphyxmini0247-transmitterlabel}
  \lean{PhyXMiniProblems.ProblemPhyXMini0247.TransmitterLabel}
  The two radio sources, distinguished by their position in the figure
\end{definition}
\begin{proof}
  This is the declared model definition of Transmitter Label.
\end{proof}
\begin{definition}[Runway Side]
  \label{def:physics:phyx-mini-0247:phyxminiproblems-problemphyxmini0247-runwayside}
  \lean{PhyXMiniProblems.ProblemPhyXMini0247.RunwaySide}
  The two sides of the runway's nodal centerline
\end{definition}
\begin{proof}
  This is the declared model definition of Runway Side.
\end{proof}
\begin{definition}[Figure Path Label]
  \label{def:physics:phyx-mini-0247:phyxminiproblems-problemphyxmini0247-figurepathlabel}
  \lean{PhyXMiniProblems.ProblemPhyXMini0247.FigurePathLabel}
  The source-to-P path labels printed in the primary figure
\end{definition}
\begin{proof}
  This is the declared model definition of Figure Path Label.
\end{proof}
\begin{definition}[Emitter Kind]
  \label{def:physics:phyx-mini-0247:phyxminiproblems-problemphyxmini0247-emitterkind}
  \lean{PhyXMiniProblems.ProblemPhyXMini0247.EmitterKind}
  Physical kind of each coherent emitter
\end{definition}
\begin{proof}
  This is the declared model definition of Emitter Kind.
\end{proof}
\begin{definition}[Propagation Medium]
  \label{def:physics:phyx-mini-0247:phyxminiproblems-problemphyxmini0247-propagationmedium}
  \lean{PhyXMiniProblems.ProblemPhyXMini0247.PropagationMedium}
  Homogeneous propagation medium used by the landing-aid model
\end{definition}
\begin{proof}
  This is the declared model definition of Propagation Medium.
\end{proof}
\begin{definition}[Receiver Response]
  \label{def:physics:phyx-mini-0247:phyxminiproblems-problemphyxmini0247-receiverresponse}
  \lean{PhyXMiniProblems.ProblemPhyXMini0247.ReceiverResponse}
  Response of the airplane's radio receiver
\end{definition}
\begin{proof}
  This is the declared model definition of Receiver Response.
\end{proof}
\begin{definition}[Guidance Region]
  \label{def:physics:phyx-mini-0247:phyxminiproblems-problemphyxmini0247-guidanceregion}
  \lean{PhyXMiniProblems.ProblemPhyXMini0247.GuidanceRegion}
  Operational regions distinguished by the guidance description
\end{definition}
\begin{proof}
  This is the declared model definition of Guidance Region.
\end{proof}
\begin{definition}[Path Difference Approximation]
  \label{def:physics:phyx-mini-0247:phyxminiproblems-problemphyxmini0247-pathdifferenceapproximation}
  \lean{PhyXMiniProblems.ProblemPhyXMini0247.PathDifferenceApproximation}
  Approximation used to turn the source geometry into a path difference
\end{definition}
\begin{proof}
  This is the declared model definition of Path Difference Approximation.
\end{proof}
\begin{definition}[Figure Feature]
  \label{def:physics:phyx-mini-0247:phyxminiproblems-problemphyxmini0247-figurefeature}
  \lean{PhyXMiniProblems.ProblemPhyXMini0247.FigureFeature}
  Qualitative objects and annotations visible in the primary image
\end{definition}
\begin{proof}
  This is the declared model definition of Figure Feature.
\end{proof}
\begin{definition}[Radio Landing Figure]
  \label{def:physics:phyx-mini-0247:phyxminiproblems-problemphyxmini0247-radiolandingfigure}
  \lean{PhyXMiniProblems.ProblemPhyXMini0247.RadioLandingFigure}
  \uses{def:physics:phyx-mini-0247:phyxminiproblems-problemphyxmini0247-transmitterlabel, def:physics:phyx-mini-0247:phyxminiproblems-problemphyxmini0247-runwayside, def:physics:phyx-mini-0247:phyxminiproblems-problemphyxmini0247-figurepathlabel, def:physics:phyx-mini-0247:phyxminiproblems-problemphyxmini0247-figurefeature}
  Qualitative content of the supplied schematic. Numerical quantities are fields of the physical setup rather than strings hidden in this figure record
\end{definition}
\begin{proof}
  This is the declared model definition of Radio Landing Figure.
\end{proof}
\begin{definition}[Radio Landing Interference Setup]
  \label{def:physics:phyx-mini-0247:phyxminiproblems-problemphyxmini0247-radiolandinginterferencesetup}
  \lean{PhyXMiniProblems.ProblemPhyXMini0247.RadioLandingInterferenceSetup}
  \uses{def:physics:phyx-mini-0247:phyxminiproblems-problemphyxmini0247-lengthquantity, def:physics:phyx-mini-0247:phyxminiproblems-problemphyxmini0247-frequencyquantity, def:physics:phyx-mini-0247:phyxminiproblems-problemphyxmini0247-speedquantity, def:physics:phyx-mini-0247:phyxminiproblems-problemphyxmini0247-transmitterlabel, def:physics:phyx-mini-0247:phyxminiproblems-problemphyxmini0247-runwayside, def:physics:phyx-mini-0247:phyxminiproblems-problemphyxmini0247-emitterkind, def:physics:phyx-mini-0247:phyxminiproblems-problemphyxmini0247-propagationmedium, def:physics:phyx-mini-0247:phyxminiproblems-problemphyxmini0247-receiverresponse, def:physics:phyx-mini-0247:phyxminiproblems-problemphyxmini0247-guidanceregion, def:physics:phyx-mini-0247:phyxminiproblems-problemphyxmini0247-pathdifferenceapproximation, def:physics:phyx-mini-0247:phyxminiproblems-problemphyxmini0247-radiolandingfigure}
  All independent physical quantities and observations in the design. intensityMaximumAtOffset side order describes the antinodal branch at the common offset magnitude, with order = 0 denoting the first branch away from the nodal line. It is related to phase and wavelength only by the governing interference law below. In particular, no setup field assigns the requested frequency or selects an answer choice
\end{definition}
\begin{proof}
  This is the declared model definition of Radio Landing Interference Setup.
\end{proof}
\begin{definition}[initial Phase Advance Radians]
  \label{def:physics:phyx-mini-0247:phyxminiproblems-problemphyxmini0247-initialphaseadvanceradians}
  \lean{PhyXMiniProblems.ProblemPhyXMini0247.initialPhaseAdvanceRadians}
  \uses{def:physics:phyx-mini-0247:phyxminiproblems-problemphyxmini0247-radiolandinginterferencesetup}
  The lower source is the longer-path source for the pictured upper point P. This helper records its initial phase advance over the upper source
\end{definition}
\begin{proof}
  This is the declared model definition of initial Phase Advance Radians.
\end{proof}
\begin{definition}[Matches Coherent Radio Source Description]
  \label{def:physics:phyx-mini-0247:phyxminiproblems-problemphyxmini0247-matchescoherentradiosourcedescription}
  \lean{PhyXMiniProblems.ProblemPhyXMini0247.MatchesCoherentRadioSourceDescription}
  \uses{def:physics:phyx-mini-0247:phyxminiproblems-problemphyxmini0247-radiolandinginterferencesetup}
  Coherent-source description: both objects are radio transmitters with one common frequency and wavelength. The specified phase offset is kept in the separate numerical-readout predicate below
\end{definition}
\begin{proof}
  This is the declared model definition of Matches Coherent Radio Source Description.
\end{proof}
\begin{definition}[Matches Problem And Figure Readouts]
  \label{def:physics:phyx-mini-0247:phyxminiproblems-problemphyxmini0247-matchesproblemandfigurereadouts}
  \lean{PhyXMiniProblems.ProblemPhyXMini0247.MatchesProblemAndFigureReadouts}
  \uses{def:physics:phyx-mini-0247:phyxminiproblems-problemphyxmini0247-lengthinmeters, def:physics:phyx-mini-0247:phyxminiproblems-problemphyxmini0247-lengthinkilometers, def:physics:phyx-mini-0247:phyxminiproblems-problemphyxmini0247-radiolandinginterferencesetup, def:physics:phyx-mini-0247:phyxminiproblems-problemphyxmini0247-initialphaseadvanceradians}
  Problem and figure readouts. The two sources and runway are 50 m across, P is 60 m from the centerline and 3.0 km = 3000 m downrange, and the lower transmitter initially leads the upper by pi radians. The same offset on both sides is stipulated to be the first intensity maximum. No frequency value occurs here
\end{definition}
\begin{proof}
  This is the declared model definition of Matches Problem And Figure Readouts.
\end{proof}
\begin{definition}[Matches Supplied Radio Landing Figure]
  \label{def:physics:phyx-mini-0247:phyxminiproblems-problemphyxmini0247-matchessuppliedradiolandingfigure}
  \lean{PhyXMiniProblems.ProblemPhyXMini0247.MatchesSuppliedRadioLandingFigure}
  \uses{def:physics:phyx-mini-0247:phyxminiproblems-problemphyxmini0247-lengthreadout, def:physics:phyx-mini-0247:phyxminiproblems-problemphyxmini0247-lengthinmeters, def:physics:phyx-mini-0247:phyxminiproblems-problemphyxmini0247-radiolandinginterferencesetup}
  Primary-image evidence: both sources, both path labels, the runway, P, and the nodal/antinodal annotations are present. The upper source-to-P ray is r1, the lower ray is r2, and P lies above the centerline
\end{definition}
\begin{proof}
  This is the declared model definition of Matches Supplied Radio Landing Figure.
\end{proof}
\begin{definition}[Matches Landing Guidance Description]
  \label{def:physics:phyx-mini-0247:phyxminiproblems-problemphyxmini0247-matcheslandingguidancedescription}
  \lean{PhyXMiniProblems.ProblemPhyXMini0247.MatchesLandingGuidanceDescription}
  \uses{def:physics:phyx-mini-0247:phyxminiproblems-problemphyxmini0247-radiolandinginterferencesetup}
  Operational interpretation of the interference pattern: the receiver is silent on the central node and produces a warning away from that line
\end{definition}
\begin{proof}
  This is the declared model definition of Matches Landing Guidance Description.
\end{proof}
\begin{definition}[Has Physical Radio Interference Parameters]
  \label{def:physics:phyx-mini-0247:phyxminiproblems-problemphyxmini0247-hasphysicalradiointerferenceparameters}
  \lean{PhyXMiniProblems.ProblemPhyXMini0247.HasPhysicalRadioInterferenceParameters}
  \uses{def:physics:phyx-mini-0247:phyxminiproblems-problemphyxmini0247-lengthinmeters, def:physics:phyx-mini-0247:phyxminiproblems-problemphyxmini0247-frequencyinhertz, def:physics:phyx-mini-0247:phyxminiproblems-problemphyxmini0247-speedinmeterspersecond, def:physics:phyx-mini-0247:phyxminiproblems-problemphyxmini0247-radiolandinginterferencesetup}
  Positivity and far-field nondegeneracy of the physical quantities
\end{definition}
\begin{proof}
  This is the declared model definition of Has Physical Radio Interference Parameters.
\end{proof}
\begin{definition}[Satisfies Path Difference Definition]
  \label{def:physics:phyx-mini-0247:phyxminiproblems-problemphyxmini0247-satisfiespathdifferencedefinition}
  \lean{PhyXMiniProblems.ProblemPhyXMini0247.SatisfiesPathDifferenceDefinition}
  \uses{def:physics:phyx-mini-0247:phyxminiproblems-problemphyxmini0247-lengthreadout, def:physics:phyx-mini-0247:phyxminiproblems-problemphyxmini0247-radiolandinginterferencesetup}
  The physical path-difference magnitude at upper point P is r2 - r1. The strict path ordering is recorded by the figure predicate and positivity conditions, so no absolute-value convention is needed here
\end{definition}
\begin{proof}
  This is the declared model definition of Satisfies Path Difference Definition.
\end{proof}
\begin{definition}[Satisfies Controlled Far Field Path Difference Model]
  \label{def:physics:phyx-mini-0247:phyxminiproblems-problemphyxmini0247-satisfiescontrolledfarfieldpathdifferencemodel}
  \lean{PhyXMiniProblems.ProblemPhyXMini0247.SatisfiesControlledFarFieldPathDifferenceModel}
  \uses{def:physics:phyx-mini-0247:phyxminiproblems-problemphyxmini0247-lengthreadout, def:physics:phyx-mini-0247:phyxminiproblems-problemphyxmini0247-radiolandinginterferencesetup}
  Controlled far-field two-source geometry. In a common length unit, the leading relation is delta r * L = d * y. The named squared-length remainder is bounded by that leading product times the dimensionless squared transverse ratios (d / (2 L))\textasciicircum{}2 + (y / L)\textasciicircum{}2. Thus the paraxial relation is local to the stated narrow-angle geometry and is not asserted as a globally exact identity. Neither the remainder relation nor its bound mentions the requested wavelength, frequency, or answer choice
\end{definition}
\begin{proof}
  This is the declared model definition of Satisfies Controlled Far Field Path Difference Model.
\end{proof}
\begin{definition}[Satisfies Out Of Phase Interference Criterion]
  \label{def:physics:phyx-mini-0247:phyxminiproblems-problemphyxmini0247-satisfiesoutofphaseinterferencecriterion}
  \lean{PhyXMiniProblems.ProblemPhyXMini0247.SatisfiesOutOfPhaseInterferenceCriterion}
  \uses{def:physics:phyx-mini-0247:phyxminiproblems-problemphyxmini0247-lengthreadout, def:physics:phyx-mini-0247:phyxminiproblems-problemphyxmini0247-runwayside, def:physics:phyx-mini-0247:phyxminiproblems-problemphyxmini0247-radiolandinginterferencesetup, def:physics:phyx-mini-0247:phyxminiproblems-problemphyxmini0247-initialphaseadvanceradians}
  Constructive-interference criterion for the out-of-phase pair. With the chosen orientation, propagation along the extra path cancels the longer-path source's initial phase advance. Thus maximum order n satisfies 2 pi delta-r / lambda = delta-phi-0 + 2 pi n. The criterion is quantified over both symmetric sides, every order, and every length unit; it contains no answer frequency
\end{definition}
\begin{proof}
  This is the declared model definition of Satisfies Out Of Phase Interference Criterion.
\end{proof}
\begin{definition}[Satisfies Radio Wave Speed Law]
  \label{def:physics:phyx-mini-0247:phyxminiproblems-problemphyxmini0247-satisfiesradiowavespeedlaw}
  \lean{PhyXMiniProblems.ProblemPhyXMini0247.SatisfiesRadioWaveSpeedLaw}
  \uses{def:physics:phyx-mini-0247:phyxminiproblems-problemphyxmini0247-lengthreadout, def:physics:phyx-mini-0247:phyxminiproblems-problemphyxmini0247-frequencyreadout, def:physics:phyx-mini-0247:phyxminiproblems-problemphyxmini0247-speedreadout, def:physics:phyx-mini-0247:phyxminiproblems-problemphyxmini0247-transmitterlabel, def:physics:phyx-mini-0247:phyxminiproblems-problemphyxmini0247-radiolandinginterferencesetup}
  Nondispersive electromagnetic-wave relation v = lambda f, for both transmitters and every compatible choice of length and time units
\end{definition}
\begin{proof}
  This is the declared model definition of Satisfies Radio Wave Speed Law.
\end{proof}
\begin{definition}[Uses Ambient Air Vacuum Propagation]
  \label{def:physics:phyx-mini-0247:phyxminiproblems-problemphyxmini0247-usesambientairvacuumpropagation}
  \lean{PhyXMiniProblems.ProblemPhyXMini0247.UsesAmbientAirVacuumPropagation}
  \uses{def:physics:phyx-mini-0247:phyxminiproblems-problemphyxmini0247-speedinmeterspersecond, def:physics:phyx-mini-0247:phyxminiproblems-problemphyxmini0247-vacuumspeedoflightinmeterspersecond, def:physics:phyx-mini-0247:phyxminiproblems-problemphyxmini0247-radiolandinginterferencesetup}
  Radio propagation through ambient air is approximated by vacuum propagation. The calibration is Physlib's exact SI speed of light, not the rounded target frequency and not an answer-choice readout
\end{definition}
\begin{proof}
  This is the declared model definition of Uses Ambient Air Vacuum Propagation.
\end{proof}
\begin{definition}[Answer Choice]
  \label{def:physics:phyx-mini-0247:phyxminiproblems-problemphyxmini0247-answerchoice}
  \lean{PhyXMiniProblems.ProblemPhyXMini0247.AnswerChoice}
  Labels of the four transmitter-frequency choices
\end{definition}
\begin{proof}
  This is the declared model definition of Answer Choice.
\end{proof}
\begin{definition}[displayed Frequency Megahertz]
  \label{def:physics:phyx-mini-0247:phyxminiproblems-problemphyxmini0247-displayedfrequencymegahertz}
  \lean{PhyXMiniProblems.ProblemPhyXMini0247.displayedFrequencyMegahertz}
  \uses{def:physics:phyx-mini-0247:phyxminiproblems-problemphyxmini0247-answerchoice}
  Frequency in megahertz displayed beside each answer label
\end{definition}
\begin{proof}
  This is the declared model definition of displayed Frequency Megahertz.
\end{proof}
\begin{definition}[recorded Dataset Answer]
  \label{def:physics:phyx-mini-0247:phyxminiproblems-problemphyxmini0247-recordeddatasetanswer}
  \lean{PhyXMiniProblems.ProblemPhyXMini0247.recordedDatasetAnswer}
  \uses{def:physics:phyx-mini-0247:phyxminiproblems-problemphyxmini0247-answerchoice}
  Answer label recorded by the source dataset
\end{definition}
\begin{proof}
  This is the declared model definition of recorded Dataset Answer.
\end{proof}
\begin{definition}[Rounds To Nearest Megahertz]
  \label{def:physics:phyx-mini-0247:phyxminiproblems-problemphyxmini0247-roundstonearestmegahertz}
  \lean{PhyXMiniProblems.ProblemPhyXMini0247.RoundsToNearestMegahertz}
  \uses{def:physics:phyx-mini-0247:phyxminiproblems-problemphyxmini0247-frequencyquantity, def:physics:phyx-mini-0247:phyxminiproblems-problemphyxmini0247-frequencyinmegahertz}
  A physical frequency rounds to a displayed whole-megahertz value when it is strictly within half a megahertz. This accommodates use of Physlib's exact 299792458 m/s speed rather than silently replacing it by 3.00e8 m/s
\end{definition}
\begin{proof}
  This is the declared model definition of Rounds To Nearest Megahertz.
\end{proof}
\begin{definition}[Matches Answer Choice]
  \label{def:physics:phyx-mini-0247:phyxminiproblems-problemphyxmini0247-matchesanswerchoice}
  \lean{PhyXMiniProblems.ProblemPhyXMini0247.MatchesAnswerChoice}
  \uses{def:physics:phyx-mini-0247:phyxminiproblems-problemphyxmini0247-radiolandinginterferencesetup, def:physics:phyx-mini-0247:phyxminiproblems-problemphyxmini0247-answerchoice, def:physics:phyx-mini-0247:phyxminiproblems-problemphyxmini0247-displayedfrequencymegahertz, def:physics:phyx-mini-0247:phyxminiproblems-problemphyxmini0247-roundstonearestmegahertz}
  Both coherent transmitters agree with the frequency printed for a choice
\end{definition}
\begin{proof}
  This is the declared model definition of Matches Answer Choice.
\end{proof}
\begin{definition}[Is Unique Matching Answer Choice]
  \label{def:physics:phyx-mini-0247:phyxminiproblems-problemphyxmini0247-isuniquematchinganswerchoice}
  \lean{PhyXMiniProblems.ProblemPhyXMini0247.IsUniqueMatchingAnswerChoice}
  \uses{def:physics:phyx-mini-0247:phyxminiproblems-problemphyxmini0247-radiolandinginterferencesetup, def:physics:phyx-mini-0247:phyxminiproblems-problemphyxmini0247-answerchoice, def:physics:phyx-mini-0247:phyxminiproblems-problemphyxmini0247-matchesanswerchoice}
  The selected answer is unique among the displayed frequencies. This generic predicate does not build in any particular label or numerical frequency
\end{definition}
\begin{proof}
  This is the declared model definition of Is Unique Matching Answer Choice.
\end{proof}
% --- Archon declaration topology end ---
% --- Archon physics formalization source end ---
